\chapter{Results}\label{results}
In this chapter we provide insights into the knowledge we managed to extract during the analysis of developer's prompts in developer-\acrshort{gpt} interaction. This chapter focuses on the results extracted from N-grams, topic modeling and sequence mining of the prompts. 

\section{N-grams}\label{sec:ngrams-result}
N-grams are extracted in order to provide an overview of the most common lexical patterns found within the data. Data gathered from each source was analyzed separately and the most common 1- to 4-grams were extracted. Table~\ref{table:common-ngrams} shows interesting tri- and quadri-grams extracted from various data sources. 

\begin{longtable}{|l|c|m{8cm}|}
    \hline
    Source & n & Found n-grams  \\
    \hline
        \multirow{2}{*}{GitHub commits} & 3 & shell script create, encode enclose result, avoid use sed, implement following feature  \\
        \cline{2-3}
        & 4 & script create change file, avoid use sed favor, task implement following feature \\
    \hline
        \multirow{2}{*}{GitHub issues} & 3 & th block th, valid string generation, state lookahead token, kernel patch instal \\
        \cline{2-3}
        & 4 & td th block th, string generation metadata information, metadata information visit input, community edition line module \\
    \hline
        \multirow{2}{*}{GitHub discussions} & 3 & technology feasibility stable, operational feasibility challenge, clause potentially unfair, employee desirability mixed/high, open/close code environment \\
        \cline{2-3}
        & 4 & technology feasibility stable employee, employee desirability high team \\
    \hline
        \multirow{2}{*}{GitHub pull requests} & 3 & gazebo sensor imu, plugin sensor accelerometer, deprecate pip enforce, possible replacement use  \\
        \cline{2-3}
        & 4 & coelemic cavity lumen subclassof, gazebo sensor imu plugin, sensor accelerometer parent origin, replacement use pip package, string sure want delete \\
    \hline
        \multirow{2}{*}{GitHub code files} & 3 & rdf type owl, restriction owl minqualifiedcardinality, passage book rewrite, break fourth wall\\
        \cline{2-3}
        & 4 & rdf type owl objectproperty, rdfs subpropertyof owl topobjectproperty, george r r martin, book rewrite passage remove, break fourth wall include, wall include reference reader \\
    \hline
        \multirow{2}{*}{GitHub repositories} & 3 & fill height viewport, increase understanding universe, help reduce suffering, research educational institution \\
        \cline{2-3}
        & 4 & height fill height viewport, leave position foreground scyscraper, reduce suffering increase prosperity, prosperity increase understanding universe, planet intellectual community train \\
    \hline
        \multirow{2}{*}{Hacker News} & 3 & capture position position, form line capture, tic tac toe, session user host, level level headline \\
        \cline{2-3}
        & 4 & form line capture position, x currently form line, headline insert concat star, concat star schedule progn \\
    \hline
\caption{Interesting 3- and 4-grams extracted from the data.}
\label{table:common-ngrams}
\end{longtable}

The extracted n-grams reveal distinct linguistic and thematic patterns across different data sources. Prompts from GitHub commits and pull requests use task-oriented, direct and action-driven communication typical of development workflows. The prompts often focus on modifying, automating or improving code artefacts. In contrast, GitHub discussions contain planning-related n-grams like ``clause potentially unfair'' and ``technology feasibility stable employee'', suggesting that contributors debate the feasibility and desirability of the solution, not just technical correctness. 

GitHub repositories show recurring themes in UI rendering (e.g.m CSS and Viewport references) as well as humanitarian, educational, and scientific discussions. However, prompts from GitHub issues and code files are influenced by embedded code or markup, such as HTML and RDF/OWL artefacts, indicating that developers often seek help with specific technical problems.

Hacker News conversations differ in a less technical, more playful approach. N-grams like ``tic tac toe'', ``form line capture position'', and ``headline insert concat star'' combine elements of game logic, syntax play, and UI behavior, typical of exploratory or playful community posts.

In conclusion, each platform reflects its users' intent, ranging from practical coding to creative exploration.

\section{Sequential pattern mining}\label{sec:spm-results}
\Acrfull{spm} is used to extend sequential pattern analysis of the data present in the datasets. For this research we consider sequences that are longer than 4 tokens (1-, 2-, 3- and 4-grams were analyzed in Section~\ref{sec:ngrams-result}) and appear in 2 or more percent of the non-empty prompts. 

These requirements were applied to the data sourced from all 7 sources. However, only data sourced from GitHub commits and repositories contained sequences that corresponded to the requirements. The data sourced from GitHub commits, issues, code files, pull requests and HackerNews contained no such sequences. Thus, in this part of the research we only focus on the data from the two sources and their contents. 

Table~\ref{table:sequences} shows categories with the most informative and recurrent sequence patterns extracted from GitHub commits and repository content. The results are grouped by source and functional themes.

\begin{longtable}{|m{15mm}|m{18mm}|m{8cm}|}
    \hline
    Source & Group & Found sequences  \\
    \hline
        & Top-5 & format shell script create everything (660), format shell script create change file everything (656), output format shell script create change file (656), format script create change file everything (656), format shell create change file everything (656) \\
        \cline{2-3}
        & Scripting automation & enclose result shell script create file solve task (645), format encode enclose result shell script create change file everything (645), shell script create change file everything solve task file small (605), export single output format encode result shell script create (121) \\
        \cline{2-3}
        GitHub commits & Constraints and guidance & avoid use sed favor full file use prevent substitution write (250), avoid use sed full file solution write concise (119) \\
        \cline{2-3}
        & Feature implementation & implement feature create new file need file (78), implement following feature create file need file (78), project specific file export single function (78) \\
        \cline{2-3}
        & Platform dependencies & task file small full file osx (78), substitution osx npm jq write text outside script (77) \\
    \hline
        & Top-5 & fill height viewport leave position foreground left edge right (10), height fill height viewport leave position foreground left edge (10), height viewport leave position foreground left edge right position (10), viewport leave position foreground left edge right position foreground (10), fill height leave position foreground left edge right (10) \\
        \cline{2-3}
        GitHub repositories & Web UI layout composition & fill height viewport leave position foreground left edge right (10), fill height viewport leave position foreground skyscraper left edge right (9), foreground skyscraper left edge right position foreground skyscraper right edge (9), landscape foundation height set height fill viewport skyscraper height fill (9) \\
        \cline{2-3}
        & Continuous research benefits & continue research work help reduce suffering planet intellectual community restore (8), disaster preparedness program action help reduce suffering increase prosperity increase (8), planet share knowledge expertise planet scientist scholar help fill knowledge (8) \\
        \cline{2-3}
        & DOM structure & stylesheet body div landscape div sun div city skyscraper div (3), print data get document embed id metadata data document data (3), embed id metadata document embed convert list (3) \\
    \hline
\caption{Interesting sequences found in the data sourced from GitHub commits and repositories.}
\label{table:sequences}
\end{longtable}

The sequence data from conversations sourced from GitHub commits and repositories shows a separation in focus: commits envode more procedure-related requests, while repositories focus more on structural and layout information. 

The most frequent patterns in GitHub commit-related conversations focus on shell scripting for file manipulation and task execution, and its constraints. Repetitive sequences like \textit{``enclose result shell script create file solve task''} and \textit{``format encode enclose result shell script create change file everything''} indicate a multistep automation pipeline for file manipulation, including file formatting, encoding, and task resolution. Additionally, commit data provides constraints aimed at maintainability, with sequences discouraging tools like \textbf{sed} in favor of full-file edits. Other patterns describe feature development practices, such as modular file creation, and even include platform-specific workflows using tools like \textbf{npm}, \textbf{jq}, or referencing \textbf{osx}.

In contrast, conversations sourced from GitHub repositories contained sequences that are dominated by UI layout and rendering logic. Patterns like ``fill height viewport leave position foreground left edge right'' reflect deterministic spatial composition, likely connected to DOM or canvas rendering engines. References to ``skyscraper'' --- a vertical stacking model, and patterns containing explicit DOM hierarchies, reinforce a component-based layout approach.

Overall, conversations from commits capture the execution trace of development, like scripts, tasks, and automation, while conversations from repositories reflect the architectural blueprint of the system, including layout logic and project narrative.

\section{Topic modelling}
To investigate how developers interact with ChatGPT across different development contexts, we used topic modeling on conversations extracted from seven sources, including GitHub-based sources (commits, issues, discussions, pull requests, code files, and repositories) and HackerNews \linebreak threads. \Acrfull{lda} topic modeling was used in order to analyze the data and extract topics covered in the data. The optimal number of topics for each part of the dataset was determined by selecting the number of topics with high coherence and low perplexity scores. Section~\ref{sec:topic-modelling-methodology} covered this process in more detail. 

Appendices~\ref{appendix:topic-models} and~\ref{appendix:topic-models-wordclouds} provide a full overview of the topic models in table and WordCloud form, respectively. The resulting topic models, summarized in Tables~\ref{table:commits-topics} to~\ref{table:hn-topics}, reveal a distribution of themes aligned with different stages of the software development lifecycle.

The topics derived from commit-based conversations are focused on actionable programming tasks, like file handling and project configuration. The \textbf{Task-based file management} topic includes words like \textit{``file''}, \textit{``task''}, \textit{``solve''}, and \textit{``create''}, which point to developer requests aimed at automating or scripting low-level actions. Similarly, \textbf{Function testing} and \textbf{UI layout planning} reflect frequent use cases in which developers ask \acrshort{gpt} for code to validate behavior or generate visual components. The \textbf{AI-assisted development} topic suggests that developers outsource code scaffolding to \acrshort{gpt}, especially for feature generation and scripting tasks.

In issue-oriented conversations, developers focus more on troubleshooting behavior and reasoning about data. Topics such as \textbf{Code \& model behavior} and \textbf{Data \& value processing} contain terms like \textit{``error''}, \textit{``input''}, \textit{``model''}, and \textit{``value''}, showing that developers use \acrshort{gpt} for clarification and debugging assistance when encountering code errors or inconsistencies. These topics underscore \acrshort{gpt}'s use as a debugging companion and computational reasoning assistant.

Discussions are broader in scope and, as noted before, have topics with a low coherence score. The \textbf{Product design feasibility} topic includes terms such as \textit{``feasibility''}, \textit{``desirability''}, and \textit{``technology''}, indicating higher-level conversations about user needs or system requirements. Several topics contain references to the \textit{Silicon Valley} series\footnote{\url{https://www.imdb.com/title/tt2575988/}}, mentioning characters and companies' names. Technical topics like \textbf{Node installation issues} and \textbf{Model calibration} show that developers also engage \acrshort{gpt} in discussions about environment setup and machine learning workflows, respectively. These themes blend the topics of software engineering, storytelling, and machine learning in discussion forums.

Conversations sourced from pull requests are described with technical precision, and deal heavily with code correctness, dependency resolution, and language semantics. Topics such as \textbf{C++ function signatures}, \textbf{Ontology \& classification}, and \textbf{Mathematical constants \& types} suggest that developers use \acrshort{gpt} to verify and generate low-level code. Other topics like \textbf{Web attestation APIs}, \textbf{Node package metadata}, and \textbf{Python packaging errors} reflect common questions about system integration, library usage, and environment misconfiguration. These patterns show how developers rely on \acrshort{gpt} for help in aligning implementation details with expected behavior.

Code file-based conversations reflect on documentation, algorithmic generation, and knowledge encoding. The \textbf{Step-by-step coding guidance} topic (e.g., \textit{``step''}, \textit{``please''}, \textit{``solution''}) shows that many conversations include instructional sequences. The presence of \textbf{OWL/RDF ontologies} suggests \acrshort{gpt} use in data modeling, and \textbf{Functional file operations} shows traditional programming tasks assistance. Meanwhile, \textbf{Evidence-supported argument} and \textbf{Book passage creation} indicate mixed-domain usage, showing that \acrshort{gpt} also supports non-code-related writing tasks within GitHub codebases. 

Topics from repositories' conversations focus on structured layouts and project-level semantics. The \textbf{UI layout metadata} topic includes tokens such as \textit{``viewport''}, \textit{``skyscraper''}, and \textit{``fill''}, reflecting UI generation and interface rendering. A second topic, \textbf{Research effort}, contains mission-driven language (e.g., \textit{``planet''}, \textit{``disaster''}, \textit{``reduce''}), indicating that repository-level content includes vision statements or long-term goals.

Topics extracted from HackerNews conversations reveal a more dialogic tone compared to other platforms. \textbf{Conversational requests} and \textbf{Narrative creation} show how developers use \acrshort{gpt} to simulate discussions and tell stories. Others, like \textbf{Game interaction errors}, \textbf{File tracking utilities} and \textbf{Code support}, return to utility-based interactions. Topics like \textbf{Company schedule} and \textbf{Positional pattern capture} contain many repetitive sequences, making them uninformative. These sourced conversations appear to blend technical querying with user-driven storytelling, positioning \acrshort{gpt} as both an assistant and a conversational partner in exploratory thinking.

Across all seven data sources, topic modeling reveals that developer-\acrshort{gpt} interactions are highly context-sensitive, aligning closely with the purpose and communication norms of each platform. While conversations sourced from GitHub commits and issues emphasize low-level execution and debugging, conversations sourced from pull requests and code files highlight correctness and structure. In contrast, \acrshort{gpt} conversations found in GitHub discussions and repositories tend to encode project-level reasoning and documentation, and HackerNews reflects informal, exploratory, or social discussions. Collectively, these patterns underscore \acrshort{gpt}'s role as a multifunctional collaborator: serving as a code assistant, design partner, documentation helper, and narrative generator.